\section{Zdroje a AI}

\vspace{1.5em}
{ \large \sloppy
\noindent\textbf{Zdroje:}

\vspace{0.75em}
\noindent [1] IAL - Algoritmy: 3. prednáška (Abstraktné dátové typy II.) Brno: Fakulta informačních technológií VUT v Brně, 2026 [cit. 2026-11-21]

\vspace{0.75em}
\noindent [2] IAL - Algoritmy: 4. prednáška (Stromové dátové štruktúry) Brno: Fakulta informačních technológií VUT v Brně, 2026 [cit. 2026-11-21]

\vspace{0.75em}
\noindent [3] IAL - Algoritmy: 6. prednáška (Vyhľadávacie tabulky II.) Brno: Fakulta informačních technológií VUT v Brně, 2026 [cit. 2026-11-21]

\vspace{0.75em}
\noindent Všetky materiály sú dostupné z e-learningového portálu VUT v Brně (Moodle), kurz Algoritmy (IAL 26/27Z)

\vspace{1.5em}

\noindent\textbf{Využitie AI:}

\vspace{0.75em}
\noindent Nástroje AI sme využívali v súlade zo zadaním. Primárne na objasnenie zadania a pri hladaní problémov a chýb.

\vspace{0.75em}
\noindent Príklad: Ahoj mám v AST navrženou strukturu pro deklaraci proměnné a teď vytvářím definici funkce, napadlo mě, že [...]. Je to správný přístup? Negeneruj kód, pouze diskutuj a klaď řečnické otázky které prohloubí mé přemýšlení o formálních jazycích. 
}